% GENERATED FILE. Edit canonical lessons, then run npm run generate:books.
% canonical-source-hash: fnv1a64:3130b5df77d8aa0a
% canonical-lessons: IT-C490-insonnia, IT-C490-russare, IT-C490-incubo, IT-C490-sbadigliare, IT-C490-sgranchirsi

\chapter{Notti agitate}
\label{ch:it-notti-agitate}
\hlchaptermodality{490}

\begin{chapteropening}
\textbf{By the end of this chapter:} \emph{I can talk about sleeping badly: insomnia, snoring and nightmares.}

Complete the last lesson of chapter 490: I can talk about sleeping badly: insomnia, snoring and nightmares.
\end{chapteropening}

\section[l'insonnia]{l'insonnia — insomnia, sleeplessness}

\label{lesson:IT-C490-insonnia}



\begin{quote}
Before the new one: say the Italian for infection, then the Italian for to recover.
\end{quote}

\subsection*{You'll want to know: l'insonnia}
\textbf{l'insonnia} — \textquotedblleft{}insomnia, sleeplessness\textquotedblright{}.

\begin{quote}
Quando sono stressato soffro d'insonnia.
\end{quote}

\begin{quote}
\emph{When I am stressed I suffer from insomnia.}
\end{quote}

\emph{Soffrire d'insonnia} — to suffer from insomnia.

\begin{cousinweb}
\emph{insonnia} is a learned borrowing from Latin \emph{insomnia} \textquotedblleft{}sleeplessness\textquotedblright{} (EtymDB), from \emph{in-} \textquotedblleft{}not\textquotedblright{} and \emph{somnus} \textquotedblleft{}sleep\textquotedblright{}.
\end{cousinweb}

\begin{grammarlens}[title={What you've built}]
One more word for telling, explaining and describing at B1.
\end{grammarlens}

\subsection*{Your turn}
\begin{itemize}
\raggedright
  \item \emph{Say it:} \emph{l'insonnia}
  \item \emph{Say it:} \emph{l'insonnia}, once more
  \item \emph{Say it:} a sentence of your own with \emph{l'insonnia}
  \item \emph{Recall:} say the Italian for infection, then the Italian for to recover, then say \emph{l'insonnia} again
\end{itemize}

\begin{tcolorbox}[breakable,colback=teal!4,colframe=teal!35!black,title={Before you move on}]
What does l'insonnia mean? (\textquotedblleft{}Insomnia, sleeplessness\textquotedblright{}.) Say it once more.
\end{tcolorbox}
\section[russare]{russare — to snore}

\label{lesson:IT-C490-russare}



\begin{quote}
Before the new one: say the Italian for to recover, then the Italian for insomnia.
\end{quote}

\subsection*{You'll want to know: russare}
\textbf{russare} — \textquotedblleft{}to snore\textquotedblright{}.

\begin{quote}
Mio marito russa e io non dormo.
\end{quote}

\begin{quote}
\emph{My husband snores and I cannot sleep.}
\end{quote}

An \emph{-are} verb: \emph{russo}, \emph{russi}, \emph{russa}.

\begin{cousinweb}
The origin of \emph{russare} is not settled; it is usually taken to imitate the sound of snoring.
\end{cousinweb}

\begin{grammarlens}[title={What you've built}]
One more word for telling, explaining and describing at B1.
\end{grammarlens}

\subsection*{Your turn}
\begin{itemize}
\raggedright
  \item \emph{Say it:} \emph{russare}
  \item \emph{Say it:} \emph{russare}, once more
  \item \emph{Say it:} a sentence of your own with \emph{russare}
  \item \emph{Recall:} say the Italian for to recover, then the Italian for insomnia, then say \emph{russare} again
\end{itemize}

\begin{tcolorbox}[breakable,colback=teal!4,colframe=teal!35!black,title={Before you move on}]
What does russare mean? (\textquotedblleft{}To snore\textquotedblright{}.) Say it once more.
\end{tcolorbox}
\section[l'incubo]{l'incubo — nightmare}

\label{lesson:IT-C490-incubo}



\begin{quote}
Before the new one: say the Italian for insomnia, then the Italian for to snore.
\end{quote}

\subsection*{You'll want to know: l'incubo}
\textbf{l'incubo} — \textquotedblleft{}nightmare\textquotedblright{}.

\begin{quote}
Stanotte ho fatto un incubo terribile.
\end{quote}

\begin{quote}
\emph{Last night I had a terrible nightmare.}
\end{quote}

\emph{Fare un incubo} — to have a nightmare. Also a bad situation: \emph{è stato un incubo}.

\begin{cousinweb}
\emph{incubo} comes from Late Latin \emph{incubus} \textquotedblleft{}a spirit that lies on the sleeper\textquotedblright{} (EtymDB), from \emph{incubāre} \textquotedblleft{}to lie upon\textquotedblright{} (Lewis \& Short).
\end{cousinweb}

\begin{grammarlens}[title={What you've built}]
One more word for telling, explaining and describing at B1.
\end{grammarlens}

\subsection*{Your turn}
\begin{itemize}
\raggedright
  \item \emph{Say it:} \emph{l'incubo}
  \item \emph{Say it:} \emph{l'incubo}, once more
  \item \emph{Say it:} a sentence of your own with \emph{l'incubo}
  \item \emph{Recall:} say the Italian for insomnia, then the Italian for to snore, then say \emph{l'incubo} again
\end{itemize}

\begin{tcolorbox}[breakable,colback=teal!4,colframe=teal!35!black,title={Before you move on}]
What does l'incubo mean? (\textquotedblleft{}Nightmare\textquotedblright{}.) Say it once more.
\end{tcolorbox}
\section[sbadigliare]{sbadigliare — to yawn}

\label{lesson:IT-C490-sbadigliare}



\begin{quote}
Before the new one: say the Italian for to snore, then the Italian for nightmare.
\end{quote}

\subsection*{You'll want to know: sbadigliare}
\textbf{sbadigliare} — \textquotedblleft{}to yawn\textquotedblright{}.

\begin{quote}
Lo spettacolo era così noioso che tutti sbadigliavano.
\end{quote}

\begin{quote}
\emph{The show was so boring that everyone was yawning.}
\end{quote}

The noun is \emph{lo sbadiglio}.

\begin{cousinweb}
\emph{sbadigliare} is usually traced to a Vulgar Latin verb \emph{bataclāre} \textquotedblleft{}to gape\textquotedblright{}, with \emph{s-} in front, a form usually linked to \emph{batāre} \textquotedblleft{}to gape\textquotedblright{}; EtymDB lists it beside French \emph{bâiller} and Occitan \emph{badalhar}, which come from \emph{bataclāre}.
\end{cousinweb}

\begin{grammarlens}[title={What you've built}]
One more word for telling, explaining and describing at B1.
\end{grammarlens}

\subsection*{Your turn}
\begin{itemize}
\raggedright
  \item \emph{Say it:} \emph{sbadigliare}
  \item \emph{Say it:} \emph{sbadigliare}, once more
  \item \emph{Say it:} a sentence of your own with \emph{sbadigliare}
  \item \emph{Recall:} say the Italian for to snore, then the Italian for nightmare, then say \emph{sbadigliare} again
\end{itemize}

\begin{tcolorbox}[breakable,colback=teal!4,colframe=teal!35!black,title={Before you move on}]
What does sbadigliare mean? (\textquotedblleft{}To yawn\textquotedblright{}.) Say it once more.
\end{tcolorbox}
\section[sgranchirsi]{sgranchirsi — to stretch (one's legs)}

\label{lesson:IT-C490-sgranchirsi}



\begin{quote}
Before the new one: say the Italian for nightmare, then the Italian for to yawn.
\end{quote}

\subsection*{You'll want to know: sgranchirsi}
\textbf{sgranchirsi} — \textquotedblleft{}to stretch (one's legs)\textquotedblright{}.

\begin{quote}
Dopo tre ore in macchina scendo a sgranchirmi le gambe.
\end{quote}

\begin{quote}
\emph{After three hours in the car I get out to stretch my legs.}
\end{quote}

\emph{Sgranchirsi le gambe} — to stretch one's legs.

\begin{cousinweb}
\emph{sgranchirsi} is built from \emph{s-} and \emph{granchio} \textquotedblleft{}crab, cramp\textquotedblright{}, usually traced to Latin \emph{cancer} \textquotedblleft{}a crab\textquotedblright{} (Lewis \& Short): to work the cramp out.
\end{cousinweb}

\begin{grammarlens}[title={What you've built}]
One more word for telling, explaining and describing at B1.
\end{grammarlens}

\subsection*{Your turn}
\begin{itemize}
\raggedright
  \item \emph{Say it:} \emph{sgranchirsi}
  \item \emph{Say it:} \emph{sgranchirsi}, once more
  \item \emph{Say it:} a sentence of your own with \emph{sgranchirsi}
  \item \emph{Recall:} say the Italian for nightmare, then the Italian for to yawn, then say \emph{sgranchirsi} again
\end{itemize}

\begin{tcolorbox}[breakable,colback=teal!4,colframe=teal!35!black,title={Before you move on}]
What does sgranchirsi mean? (\textquotedblleft{}To stretch (one's legs)\textquotedblright{}.) Say it once more.
\end{tcolorbox}
